\documentclass[aspectratio=169]{beamer}
\usetheme{Copenhagen}
\usecolortheme{whale}
\usefonttheme[onlymath]{serif}

\usepackage{setspace}
\usepackage{listings}
\usepackage{xcolor}
\usepackage{graphicx}
\usepackage{multimedia}
\usepackage{fontawesome}
\usepackage[most]{tcolorbox}
\usepackage{empheq}
\usepackage{fancybox}
\usepackage{annotate-equations}
\usepackage{tikz}
\usetikzlibrary{positioning, calc, shapes.geometric, shapes.multipart, 
	shapes, arrows.meta, arrows, 
	decorations.markings, external, trees}
\tikzstyle{Arrow} = [
thick, 
decoration={
	markings,
	mark=at position 1 with {
		\arrow[thick]{latex}
	}
}, 
shorten >= 3pt, preaction = {decorate}
]

\definecolor{teel}{HTML}{00AEB3}
\definecolor{redorange}{HTML}{F26035}
\definecolor{darkred}{HTML}{8B0000}
\definecolor{darkblue}{HTML}{00008B}
\definecolor{darkgreen}{HTML}{06402B}
\definecolor{CarolinaBlue}{HTML}{4B9CD3}
\definecolor{deepsea}{HTML}{123456}
\definecolor{sharedblue}{HTML}{009292}
\definecolor{distalblue}{HTML}{490092}
\definecolor{codegreen}{rgb}{0,0.6,0}
\definecolor{codegray}{rgb}{0.5,0.5,0.5}
\definecolor{codepurple}{rgb}{0.58,0,0.82}
\definecolor{backcolour}{rgb}{0.95,0.95,0.92}

\usecolortheme[named=CarolinaBlue]{structure}

\lstdefinestyle{mystyle}{
	language=Python,
	backgroundcolor=\color{backcolour},   
	commentstyle=\color{codegreen},
	keywordstyle=\color{magenta},
	numberstyle=\tiny\color{codegray},
	stringstyle=\color{codepurple},
	basicstyle=\ttfamily\scriptsize,
	breakatwhitespace=false,         
	breaklines=true,                 
	captionpos=b,                    
	keepspaces=true,                 
	numbers=left,                    
	numbersep=5pt,                  
	showspaces=false,                
	showstringspaces=false,
	showtabs=false,                  
	tabsize=2
}

\lstset{style=mystyle}


\title[Simulation Code]{Organizing Code for Simulation Studies}
\author[Paul Zivich]{Zivich Research Group}
\date{~}

\setbeamercovered{transparent}
\setbeamertemplate{navigation symbols}{}  % gets rid of the dumb navigation symbols
\setbeamertemplate{page number in head/foot}{\insertframenumber}  % adds slide #
\setbeamertemplate{headline}{}

\AtBeginSection[]{
	\begin{frame}
		\vfill
		\centering
		\begin{beamercolorbox}[sep=8pt,center,shadow=true,rounded=true]{title}
			\usebeamerfont{title}\insertsectionhead\par%
		\end{beamercolorbox}
		\vfill
	\end{frame}
}

\begin{document}

\maketitle

\begin{frame}{Simulation Studies}
	Simulation studies are a key aspect of methods work
	\begin{itemize}
		\item Verify expected statistical behaviors
		\item Illustrate important concepts
		\item Compare finite-sample performance between competing methods
		\item Check whether code works as intended
		\item Perform power calculations
	\end{itemize}
	~\\
	Scenario should guide choices like: data generating mechanism, performance metrics, etc.
	\begin{itemize}
		\item Once those are decided, consider implementation
	\end{itemize}
\end{frame}

\begin{frame}{Code Organization}
	An organization to consider
	~\\~\\
	\texttt{\noindent Project/ \\
		\qquad Code/ \\
		\qquad\qquad output/ \\
		\qquad\qquad data\_generating\_mechanism.file \\
		\qquad\qquad estimators.file \\
		\qquad\qquad compute\_truth.file \\
		\qquad\qquad run\_simulation\_scenarios.file \\
		\qquad\qquad analyze\_scenarios.file \\
		\qquad\qquad helper.file \\
	}
	~\\~\\
	\textbf{Philosophy}: each aspect lives in own file, so any changes carry forward everywhere
\end{frame}

\begin{frame}{Data Generating Mechanism}
	Script contains all functionality to generate simulated data
	\begin{itemize}
		\item Functions to generate truth and observed data
		\item One overall function to generate complete data set
		\begin{itemize}
			\item Basic argument to include: $n$, sample size
			\item Optional: scenario when have many options
		\end{itemize}
		\item Other functions can be included, but only the overall constructor should be imported elsewhere
	\end{itemize}
\end{frame}

\begin{frame}{Estimators}
	Script containing all estimators under consideration
	\begin{itemize}
		\item Separate functions for each estimator
		\item Generic inputs for future flexibility
	\end{itemize}
	~\\
	Helper functions
	\begin{itemize}
		\item Script that contains any helper functions you might need for a simulation
	\end{itemize}
\end{frame}

\begin{frame}{Running Simulations}
	Compute the truth
	\begin{itemize}
		\item Own script that computes the truth and saves it as a reference (like in \texttt{output})
	\end{itemize}
	~\\
	Simulation scenarios
	\begin{itemize}
		\item Exact organization depends on run-time
		\item Best to have a general script
		\begin{itemize}
			\item Script has a key variable which determines the scenario
		\end{itemize}
		\item Only calls in objects from other files
		\item Saves all simulation outputs to \texttt{output/}
	\end{itemize}
\end{frame}

\begin{frame}{Analyze Outputs}
	Separate script that
	\begin{itemize}
		\item Loads in results from \texttt{output}
		\item Creates all summary tables and figures
	\end{itemize} 
\end{frame}

\begin{frame}{Example}
	A basic (bad) example: \\
	\url{github.com/pzivich/publications-code/tree/master/ClusteringEE/simulations}
	~\\~\\
	A basic (good) example: \\
	\url{github.com/pzivich/publications-code/tree/master/PooledLogitEE/sims}
	~\\~\\
	A complex (good) example: \\
	\url{github.com/pzivich/publications-code/tree/master/CausalRNG}
\end{frame}

\begin{frame}{Questions?}
	~
\end{frame}

\end{document}
